% ============================================================
%  Chapter fragment -- included by main.tex
%  (not a standalone document; compile main.tex to build the PDF)
% ============================================================

\sheettitle{05 -- PenTesting 2: Intelligence Gathering \& Scanning}{SWS2 -- ZHAW / Tellenbach}

% ==================== INTELLIGENCE GATHERING ====================
\section{Intelligence Gathering \hl{(!)}}
PTES phase 2 (a.k.a. \kw{footprinting / reconnaissance}). Mostly \kw{OSINT} (open-source intel). Feeds the later \kw{vuln analysis} \& \kw{exploitation} phases.
\textit{Limitations:} OSINT can be \kw{inaccurate / outdated / manipulated}; some methods (e.g. dumpster diving) out of scope.

\subsection{3 Methods (by detectability) \hl{(!)}}
\begin{itemize}
\item \kw{Passive}: use \kw{3rd-party} data only, never touch the target $\to$ \kw{undetectable}. \term{Ex: Shodan, WHOIS, search engines.}
\item \kw{Semi-passive}: \kw{legitimate-looking} queries, stick to normal protocol/traffic, anonymize $\to$ should \kw{not} raise alarms. \term{Ex: normal DNS lookups.}
\item \kw{Active}: directly probe target, \kw{should be detected}. \term{Ex: vuln scanning, port scan, enumeration.}
\end{itemize}

\subsection{Intelligence Gathering Levels}
\begin{itemize}
\item \kw{L1 -- Compliance driven}: mostly \kw{automated tools}; minimal effort (a "click-button" baseline).
\item \kw{L2 -- Best practice}: automated + \kw{manual} analysis; understand the business.
\item \kw{L3 -- State sponsored}: in-depth \kw{manual}, cultivate social networks; full-scope / red-team.
\end{itemize}

\subsection{What to gather}
Layered, from outside in: \kw{Physical} (locations) $\to$ \kw{Logical} (org structure, business relations) $\to$ \kw{Financial} $\to$ \kw{Electronic documents} (metadata!) $\to$ \kw{Infrastructure assets} (domains, IPs, tech) $\to$ \kw{Employees / HUMINT}.

\subsection{Browsing the target's website}
Company site, search functions, profile/team pages $\to$ names \& roles for \kw{social engineering}. Even \kw{removed} content: \kw{Wayback Machine}, \kw{Google cache}, mementoweb.

% ==================== GOOGLE DORKING ====================
\section{Search Engine Hacking (Google Dorking)}
Craft search queries to surface \kw{exposed} files/pages. \kw{Method: Passive, Level L1/L2.}

\subsection{Operators \hl{(!)}}
{\scriptsize
\kw{site:} (limit to domain) \textbullet{} \kw{inurl:} / \kw{allinurl:} \textbullet{} \kw{intext:} / \kw{allintext:} \textbullet{} \kw{intitle:} \textbullet{} \kw{filetype:} (e.g. pdf, xls) \textbullet{} \kw{related:} \textbullet{} \kw{cache:} \textbullet{} \kw{"..."} exact \textbullet{} \kw{-} exclude \textbullet{} \kw{OR} / \kw{|} \textbullet{} \kw{*} wildcard \textbullet{} \kw{..} number range \textbullet{} \kw{AROUND(n)} proximity
}

\subsection{GHDB \& automation}
\begin{itemize}
\item \kw{Google Hacking Database (GHDB)} on \kw{exploit-db.com}: ready-made dorks by category (login portals, error messages, files with passwords...).
\item Tool: \kw{Google Hacking Diggity}.
\item \hl{Automation is hard}: throttling, \kw{CAPTCHAs} ("I'm not a robot" after suspicious traffic), query limits/cost; dorks have short \kw{lifetime}.
\end{itemize}

% ==================== INFRASTRUCTURE ASSETS ====================
\section{Infrastructure Assets}
Start from a \kw{URL or company name}. Relevant: \kw{domains}, \kw{IP addresses} \& network provider(s), external infra profile, \kw{(defence) technologies} used.

% ==================== FINDING DOMAINS ====================
\section{Finding Domains}
\hl{Problem:} no \kw{public register} of all domains. For \kw{gTLD} (.com/.net/.org) you can get \kw{zone files} (lists 2nd-level domains); for \kw{ccTLD} only a few countries publish them.
\textbf{Goal:} find all domains (= potential \kw{entry points}) owned by the company / an employee, or by \kw{B2B} partners.

\subsection{Approaches (depend on what you know)}
Company/employee name or email \textbullet{} a set of IPs \textbullet{} a domain (find subdomains) \textbullet{} a domain (find related domains).

\subsection{From company/employee $\to$ WHOIS \hl{(!)}}
\begin{itemize}
\item \kw{WHOIS} = protocol to query the \kw{registered owner / assignee} of an Internet resource (domain, IP block, AS).
\item \kw{RIRs} (Regional Internet Registries) run WHOIS servers: \kw{ARIN} (N.America), \kw{RIPE NCC} (Europe), \kw{APNIC} (Asia-Pac), \kw{LACNIC} (Latin America), \kw{AfriNIC} (Africa).
\item Entries are \kw{cross-referenced} (ARIN query for a RIPE record returns a pointer to RIPE).
\item No standard to find the \kw{responsible} server $\to$ manual: start at \kw{whois.iana.org} $\to$ query TLD (e.g. \kw{ch} $\to$ org = \kw{SWITCH}, server whois.nic.ch) $\to$ query that org's DB.
\item Query via \kw{whois} CLI or web front-ends (e.g. nic.ch/whois for .ch).
\end{itemize}

\subsection{Reverse WHOIS}
WHOIS normally: \kw{domain $\to$ owner}. Usually \kw{no official reverse lookup} (name/email $\to$ domains), and registrations may be \kw{privacy-protected} (\kw{GDPR}). $\to$ use 3rd-party: \kw{viewdns.info}, \kw{whoisxmlapi.com}. Caveat: may be \kw{incomplete} or list \kw{too many} domains.

\subsection{How useful is WHOIS? \hl{(!)}}
\begin{itemize}
\item \kw{Domain WHOIS = often thin}: since \kw{GDPR} the owner fields are \kw{redacted} / hidden behind a privacy proxy or the registrar (you see "switchplus AG", not the real holder). Still get: \kw{name servers}, registrar, dates.
\item \kw{IP / AS WHOIS = still gold}: \hl{not} redacted (GDPR protects \kw{people}, not netblocks) $\to$ querying one IP returns the \kw{whole owned range} (\kw{inetnum}) + org + AS. Bigger prize than one domain's owner field.
\end{itemize}

\subsection{From a set of IPs $\to$ DNS (PTR) \hl{(!)}}
DNS forward = \kw{name $\to$ IP} (A record). \kw{PTR record} = the \kw{reverse}, \kw{IP $\to$ hostname}, stored in the \kw{in-addr.arpa} zone (IP written backwards):
{\scriptsize
\begin{itemize}
\item \kw{A}: \texttt{www.zhaw.ch} $\to$ 160.85.104.96
\item \kw{PTR}: \texttt{96.104.85.160.in-addr.arpa} $\to$ \texttt{wwwrl.zhaw.ch}
\end{itemize}}
\hl{Not mandatory}: every live host needs an A record, but the IP owner need not publish a PTR (or it's a generic provider name, e.g. \texttt{ec2-..amazonaws.com}).

\textbf{Use in recon -- reverse-DNS sweep:} once WHOIS/BGP told you the company \kw{owns} a block, PTR-lookup \kw{every IP} in it to turn "they own this /16" into a \kw{list of named, live machines}:
{\scriptsize
\begin{itemize}
\item 160.85.104.96 $\to$ \texttt{wwwrl.zhaw.ch}
\item 160.85.104.121 $\to$ \texttt{srv-mail-011.zhaw.ch}
\item 160.85.104.245 $\to$ \texttt{srv-app-303-data2.zhaw.ch}
\end{itemize}}
Names \kw{leak the internal map}: \texttt{srv-mail-}=mail, \texttt{srv-app-}=app, \texttt{ns1/ns2}=DNS. A PTR existing $\approx$ IP is in service. Tool: \kw{robtex.com} (shows A record + all names on the same IP $\to$ related domains).

\subsection{Subdomains of a domain}
\hl{NO official way} (would need the \kw{authoritative name server}'s data). Combine \kw{unofficial} methods:
\begin{itemize}
\item \kw{Web-page scraping} (crawl main domain, grab subdomains).
\item \kw{Search engines} (design queries).
\item \kw{Brute-force} enumeration (mail., dev., ...).
\item \kw{X.509 certs}: read \kw{Subject Alternative Name (SAN)} -- one cert often lists many (sub)domains.
\end{itemize}

\subsection{Related domains}
Domains on the \kw{same IP} \textbullet{} same \kw{NS} (own NS $\to$ few \& relevant; provider's NS $\to$ many irrelevant) \textbullet{} same \kw{MX} (mailserver) \textbullet{} relations from web-page data. 3rd-party combine sources (whoisxmlapi, findsubdomains, spyse). \kw{Method: passive (WHOIS/SE/3rd-party) + active (DNS), L1/L2.}

% ==================== FINDING IP ADDRESSES ====================
\section{Finding IP Addresses}

\subsection{Company OWNS IP blocks / its own AS}
\begin{itemize}
\item IP ranges assigned by \kw{RIRs}; info includes contact + the \kw{Autonomous System (AS)} the block belongs to.
\item \kw{IP $\to$ network prefix}: RIR websites / WHOIS (e.g. RIPE query for an IP $\to$ \kw{inetnum} range like 160.85.0.0--160.85.255.255).
\item \kw{Company name $\to$ "all" IP ranges}: \kw{BGP Toolkit} (\kw{bgp.he.net}, Hurricane Electric) -- shows prefixes (v4 \& v6) per AS. Needed because enumerating IPv6 by brute force is infeasible $\to$ use \kw{BGP} routing data.
\end{itemize}

\subsection{Company does NOT own blocks/AS}
\begin{itemize}
\item \kw{Domains $\to$ IPs}: \kw{Forward-DNS} (host IP + NS/MX servers).
\item \kw{IP $\to$ more IPs}: \kw{Reverse-DNS} on the network block + other checks.
\item \kw{Search services} (by company name): \kw{Search Engine Hacking} (web-only, incomplete), \kw{Shodan / Censys}.
\end{itemize}

\subsection{DNS tools}
\begin{itemize}
\item \kw{nslookup <IP>}: reverse $\to$ in-addr.arpa $\to$ hostname.
\item \kw{dig ns <domain>}: name servers; \kw{dig mx <domain>}: mailservers (e.g. zhaw $\to$ outlook.com).
\item \kw{Reverse-DNS sweep}: script PTR lookups over a whole \kw{/16} (class B) to harvest active hostnames. \hl{Is this active or passive?} $\to$ \kw{passive} (you query \kw{DNS}, not the target's own systems).
\end{itemize}
\kw{Method: passive (RIR/3rd-party), semi-active/active (DNS), L1/L2/(L3).} Passive discovery is \kw{incomplete} $\to$ verify.

% ==================== SCANNING ====================
\section{Scanning}
Take recon results and examine networks \& hosts \kw{in detail}. Goals: determine \kw{network structure}, find \kw{reachable hosts} (depends on tester location -- inside/outside), and \kw{analyse hosts} (OS, services, software $\to$ hints at vulns).

\subsection{Network structure}
\begin{itemize}
\item Reveals interesting \kw{areas} (e.g. a \kw{DMZ}) to prioritize.
\item Provides \kw{attack paths}: a host you can't reach directly may be reached via a \kw{stepping stone} in its network.
\item \kw{traceroute}: lists all hops to target. Default \kw{UDP or ICMP} (\kw{-I}); \kw{tcptraceroute} uses TCP $\to$ try both, esp. through firewalls.
\end{itemize}
\textit{Deductions from a trace:} two hosts behind the same router are \kw{likely} same network; IPs 104.96 \& 104.121 $\Rightarrow$ at least a \kw{/25}. An \kw{asterisk} = no answer $\Rightarrow$ "\kw{some filtering}" (firewall) at that hop. Scan also \kw{internal$\to$internal} and \kw{internal$\to$external} for more structure.

\subsection{hping3}
More flexible than traceroute: pick \kw{protocol (ICMP/UDP/TCP)} \& \kw{dest port}. \term{Ex: TCP port 80 SYN traceroute} -- \kw{port 80 gets through firewalls} more often than UDP/ICMP.

% ==================== NMAP ====================
\section{nmap \hl{(!)}}

\subsection{Host discovery (ping scan) \hl{(!)}}
\kw{nmap -sn <target>} = "\kw{s}can, \kw{n}o port scan". Answers \kw{one} question: \kw{is the host alive?} -- then stops (no port table). \textit{(old syntax: -sP)}. First resolves the name to an IP, then probes it.

It doesn't send just \kw{one} ping -- it fires a \kw{bundle} and counts the host as \kw{up} if \kw{anything} replies (firewalls block some probe types, so it tries several):
{\scriptsize
\begin{itemize}
\item \kw{ICMP echo} -- the classic ping.
\item \kw{TCP SYN $\to$ 443} -- SYN/ACK back = alive.
\item \kw{TCP ACK $\to$ 80} -- RST back = alive (slips past some firewalls).
\item \kw{ICMP timestamp} -- backup when echo is blocked.
\end{itemize}}
\hl{Local network:} switches to \kw{ARP} requests (can't reach a LAN IP without its MAC; ARP can't be L3-filtered).

\textbf{Mental model:}
\begin{itemize}
\item \kw{-sn} = knock, see who's home, leave (discovery only) $\to$ then port-scan only those that answered.
\item \kw{-Pn} = skip the knock, assume \kw{everyone} is home, port-scan all. Slower, but \hl{pentest-safe}: a host that ignores all discovery probes is marked "down" \& \kw{skipped} by default $\to$ -Pn so no live host slips through.
\item \kw{default} (neither) = knock first, then scan whoever answered.
\end{itemize}

\subsection{Scan techniques \hl{(!)}}
\begin{itemize}
\item \kw{-sS} TCP SYN: \kw{default}, fast, stealthy = \kw{half-open}.
\item \kw{-sT} TCP Connect: uses OS \kw{connect()}; slower, \kw{more detectable} (no raw packets / no root).
\item \kw{-sA} TCP ACK: test \kw{firewall rulesets}.
\item \kw{-sU} UDP: much slower than TCP.
\item \kw{-sN/-sF/-sX} Null/FIN/XMAS: odd flag combos to slip past \kw{stateless} firewalls.
\item \kw{-sO} IP protocol: which IP protocols the target supports.
\end{itemize}

\subsection{Port discovery}
\kw{IANA ranges}: well-known 1--1023, registered 1024--49151, dynamic 49152--65535. Selection: \kw{-p-} (all), \kw{-p22,53}, \kw{-p20-25}, default = \kw{top 1000}, \kw{-F} = top 100, \kw{--top-ports 10}. Combine: \kw{-sU -sS -p U:53,T:21,80}.
\textit{Effectiveness:} top-1000 TCP $\approx$93\%, top-1017 UDP $\approx$100\% $\to$ \kw{fewer ports} already catch most services.

\subsection{Service / version detection}
\kw{-sV} (in \kw{-A}): interrogate open ports with \kw{service probes}. Uses the \kw{nmap-service-probes} DB of probes + \kw{match expressions}; mostly parses the \kw{banner} (connect response). Plain port label can lie (apps on arbitrary ports) $\to$ \kw{-sV} gives real product+version. Tip: \kw{--reason} = why a port is open/closed/filtered.

\subsection{OS detection}
\kw{-O -v}: sends up to \kw{16 TCP/UDP/ICMP probes} exploiting \kw{RFC ambiguities}; combines response attributes into an OS \kw{fingerprint}.

\subsection{Get it all: -A}
\kw{-A} = OS detection + version detection + \kw{script scanning} + traceroute. Yields protocols, app names, versions, hostname, device type, OS family, \kw{CPE}.

\subsection{Manual banner grabbing}
Also good for versions: \kw{nc}/\kw{telnet} (e.g. ssh banner "SSH-2.0-OpenSSH..."), \kw{openssl s\_client} for TLS. \term{Ex: server header "BigIP" $\Rightarrow$ a \kw{WAF} (F5) sits in front.}

\subsection{Passive scanning}
\kw{Shodan} (shodan.io) -- search engine for \kw{Internet-connected devices}: scans the net, grabs \kw{banners}; historic data; free tier limited. \kw{Censys} (censys.io) -- similar, more \kw{frequently updated}, only \kw{current} data. Both queried by net/company $\to$ hosts, ports, products, certs (SAN names).

% ==================== FOOTPRINTING DEFENCES ====================
\section{Footprinting Defences}
Find the \kw{defensive systems} (more attack vectors + know how to \kw{evade detection}): firewalls, \kw{WAF}, IDS, anti-virus.
\begin{itemize}
\item Brute-force protections, (D)DoS misuse, auth type/protocol.
\item \kw{MFA} $\Rightarrow$ attack vector must also obtain the \kw{2nd factor}.
\item \kw{Federated} systems $\Rightarrow$ attack via other \kw{members of the federation}.
\item \kw{Active}: inspect banners/HTTP responses, device fingerprinting, social engineering. \kw{Passive}: search-engine hacking, Censys/Shodan.
\end{itemize}
\term{Ex WAF:} manual = inspect HTTP response (Server: BigIP); automated = \kw{nmap http-waf-fingerprint} ("Detected WAF: F5 BigIP").

% ==================== HUMINT ====================
\section{HUMINT}
Human intelligence = \kw{direct interaction}: \kw{physical} (observation) or \kw{verbal} (call/chat/on-site). First learn all you can about a person: \kw{online accounts}, social-media posts.
\textbf{Automate account discovery}: sites that tell if a \kw{username exists} on a service (the service must \kw{leak} existence -- many do via login or \kw{password-reset}). Tools: \kw{knowem.com}, \kw{checkusernames.com}.

\begin{defbox}
\textbf{Take-away:} Intelligence gathering is the \kw{first} pentest step. After it you know domains, IP ranges, technical contacts, hostnames/IPs -- and sometimes critical info \kw{voluntarily leaked} by employees. Scanning then maps the network \& hosts (OS, services, versions, likely vulns) using \kw{traceroute, nmap, Shodan/Censys}.
\end{defbox}

